% Einheit 2 von 4 – Sinus- und Kosinusfunktion und ihre Merkmale (bank/trigonometrische-funktionen/e2.jsonl, zusammenbau v0.1)
%% TODO Zweigzeile Teil 1: was hier gelernt wird (Satz in Schülersprache aus dem Einheitstitel) – Einheit 2
\einheitenkopf[e2]{Einheit 2 von 4 · Sinus- und Kosinusfunktion und ihre Merkmale}
\zweigzeile{ab Kl. 10 · keine P10-Aufgabe}
\setcounter{aufgabe}{15}

%% TODO Titel: Ich-kann-Satz fehlt in der Bank (Platzhalter: Kettenname) – Nr. 16
%% TODO Anweisung: ein Satz über der ersten Teilaufgabe fehlt in der Bank – Nr. 16
\begin{aufgabe}{Sinus- und Kosinusfunktion}
\begin{teile}
\teil Wo ist die Welle zu Ende? Welcher Abschnitt ist genau eine ganze Wiederholung?\\ \kreuz{von A bis B}\\ \kreuz{von A bis C}
\end{teile}
\swa{Fülle die Wertetabelle für y = sin(x) aus (zwei Stellen).}{\wertetabelle{x in °}{y}{0,45,90,135,180,225,270,315,360}}
\swa{Fülle die Wertetabelle für y = cos(x) aus (zwei Stellen).}{\wertetabelle{x in °}{y}{0,45,90,135,180,225,270,315,360}}
\swa{Fülle die Wertetabelle für y = sin(x) aus (zwei Stellen).}{\wertetabelle{x in °}{y}{0,60,120,180,240,300,360}}
\swa{Fülle die Wertetabelle für y = cos(x) aus (zwei Stellen).}{\wertetabelle{x in °}{y}{0,60,120,180,240,300,360}}
\swa{Fülle die Wertetabelle für y = sin(x) aus (zwei Stellen).}{\wertetabelle{x in °}{y}{0,40,80,120,160,200,240,280,320,360}}
\swfrage{Was bleibt gleich, was ändert sich?}
\swa{Trage die Punkte der Tabelle ein und zeichne die Welle y = sin(x).}{\wertetabelle[0,0.71,1,0.71,0,-0.71,-1,-0.71,0]{x in °}{y}{0,45,90,135,180,225,270,315,360} \begin{ksys}[xmin=0,xmax=360,ymin=-1.2,ymax=1.2,xstep=45,ystep=0.2]\end{ksys}}
%% TODO Darstellung für schwach fehlt in der Bank (trigonometrische-funktionen-e2-k1-s3-v1; Abschnitt „Für schwache Schüler“)
\swz{Rechtsachse für einen ganzen Vollkreis ab 0°, 9 cm Platz – wie viele Grad je Zentimeter? Welche Zahlen schreibst du an?}{}
\swz{Lies sin 135° am Graphen ab (eine Stelle).}{\begin{ksys}[xmin=0,xmax=360,ymin=-1.2,ymax=1.2,xstep=45,ystep=0.2,ablesen]\funktion{sin(\x)}{f}\end{ksys}}
\swz{Lies ab: Bei welchen Winkeln im ersten Vollkreis ist sin x ≈ 0,7?}{\begin{ksys}[xmin=0,xmax=360,ymin=-1.2,ymax=1.2,xstep=45,ystep=0.2,ablesen]\funktion{sin(\x)}{f}\end{ksys}}
%% TODO Darstellung für schwach fehlt in der Bank (trigonometrische-funktionen-e2-k1-s6-v1; Abschnitt „Für schwache Schüler“)
\swz{Liegt P(50|0,77) auf dem Graphen von y = sin(x)? Rechne im Grad-Modus.}{}
%% TODO Darstellung für schwach fehlt in der Bank (trigonometrische-funktionen-e2-k1-s7-v1; Abschnitt „Für schwache Schüler“)
\swz{Welche Werte kann y = sin(x) annehmen?}{}
%% TODO Darstellung für schwach fehlt in der Bank (trigonometrische-funktionen-e2-k1-s8-v1; Abschnitt „Für schwache Schüler“)
\swz{Gib alle Nullstellen von y = sin(x) zwischen 0° und 720° an. Welche Regelmäßigkeit siehst du?}{}
%% TODO Darstellung für schwach fehlt in der Bank (trigonometrische-funktionen-e2-k1-s9-v1; Abschnitt „Für schwache Schüler“)
\swz{Gib Hoch- und Tiefpunkt von y = sin(x) im ersten Vollkreis ab 0° an.}{}
%% TODO Darstellung für schwach fehlt in der Bank (trigonometrische-funktionen-e2-k1-s10-v1; Abschnitt „Für schwache Schüler“)
\swz{Nach welchem Winkel wiederholt sich y = sin(x) zum ersten Mal ganz?}{}
\begin{teile}
\teil y = sin(x) zwischen 100° und 250° – steigt oder fällt die Kurve?\\ \kreuz{steigt}\kreuz{fällt}
\end{teile}
%% TODO Darstellung für schwach fehlt in der Bank (trigonometrische-funktionen-e2-k1-s12-v1; Abschnitt „Für schwache Schüler“)
\swz[3]{Beschreibe die Symmetrie der Sinuskurve und belege sie mit einem Wertepaar.}{}
\swz{Um wie viel Grad muss man die Sinuskurve nach links schieben, damit sie auf der Kosinuskurve liegt?}{\begin{ksys}[xmin=0,xmax=360,ymin=-1.2,ymax=1.2,xstep=45,ystep=0.2,ablesen]\funktion{sin(\x)}{f}\funktion{cos(\x)}{g}\end{ksys}}
\begin{teile}
\teil Welche Gleichung gehört zum Graphen?\\ \kreuz{y = sin(x)}\\ \kreuz{y = cos(x)}
\end{teile}
\swa{Fülle die Wertetabelle für y = sin(x) aus (zwei Stellen), zeichne den Graphen und gib Wertebereich, Nullstellen, Hoch- und Tiefpunkt, kleinste Periode und Symmetrie an.}{\wertetabelle{x in °}{y}{-180,-135,-90,-45,0,45,90,135,180} \begin{ksys}[xmin=-180,xmax=180,ymin=-1.2,ymax=1.2,xstep=45,ystep=0.2]\end{ksys}}
\end{aufgabe}

%% TODO Titel: Ich-kann-Satz fehlt in der Bank (Platzhalter: Kettenname) – Nr. 17
%% TODO Anweisung: ein Satz über der ersten Teilaufgabe fehlt in der Bank – Nr. 17
\begin{aufgabe}{Sinus- und Kosinusfunktion – Fehler finden, Begründen, Darstellungswechsel}
\swz[3]{Mehmet zeichnet die Sinuskurve so, dass sie im Hochpunkt spitz zuläuft wie ein Dach. Finde den Fehler.}{}
\swfrage{Begründe, warum sich die Sinuskurve nach einem Vollkreis wiederholt.}
\begin{teile}
\teil Gehört die Wertetabelle zu f oder zu g?\\ \kreuz{f}\kreuz{g}
\end{teile}
\end{aufgabe}

\uebersichtskasten{Sinusfunktion: y = sin(x) ordnet jedem Winkel seinen Sinuswert zu. Der Graph ist eine Welle, die zwischen −1 und 1 hin und her schwingt und sich nach einem Vollkreis wiederholt – die kleinste Periode ist 360°. \\ sin 0° = 0 \quad sin 90° = 1 \quad sin 180° = 0 \quad sin 270° = −1 \quad sin 360° = 0 \\ Kosinusfunktion: y = cos(x) hat dieselbe Wellenform, beginnt aber oben. \\ cos 0° = 1 \quad cos 90° = 0 \quad cos 180° = −1 \\ Merkmale: Man darf jeden Winkel einsetzen; die Werte liegen immer zwischen −1 und 1; die Sinuskurve hat ihre Nullstellen alle 180° und ist punktsymmetrisch zum Ursprung, die Kosinuskurve ist achsensymmetrisch zur senkrechten Achse.}
